\begin{abstract}
Jailbreak research increasingly relies on LLM judges to determine whether a
model response fulfills a malicious goal. Yet the same response can receive
different verdicts across judges, making reported attack success sensitive to
the evaluator. We ask whether such disagreements can be turned into an
automatic signal for discovering judge vulnerabilities. We introduce
\textsc{TVD}-Conflict, an agentic framework that searches for realistic
query--response pairs on which two judges, operating under the same evaluation
rubric, disagree. Their judgments provide iterative feedback for revising each
candidate, allowing the framework to collect difficult cases without requiring
human labels during the search. Experiments across ten jailbreak attacks, four
victim models, three judge backbones, and two established evaluation protocols
show that evaluator choice can substantially alter measured attack success.
Across 50 malicious goals, \textsc{TVD}-Conflict discovers 1,702 conflict
pairs that expose recurring vulnerabilities in how judges interpret rubrics,
respond to stylistic and refusal cues, reason about obfuscated content, and
ground their verdicts in evidence. These findings show that disagreement is not
merely evaluation noise, but a scalable signal for constructing diagnostic
stress tests and improving the reliability of jailbreak evaluation.
\end{abstract}
